% abstract.tex -- reframed for Applied Mathematics PhD submission
\begin{abstract}

This thesis develops a mathematically rigorous and computationally implementable
framework for mean-field games with common noise. The central object of study is the
\emph{stochastic mean-field equilibrium}: the Nash equilibrium of a large population
of interacting agents subject to both idiosyncratic and common sources of randomness,
characterised through a system of forward-backward stochastic differential equations
(FBSDEs) coupled with a stochastic Fokker--Planck equation. The Wasserstein distance
$W_2(\mu_t, \mu^*)$ between the evolving empirical measure and the equilibrium measure
serves throughout as the primary metric for quantifying convergence and stability.

The thesis establishes five principal contributions:

\begin{enumerate}
    \item \textbf{Quantitative well-posedness} (Chapter 3, Theorem 3.1).
    Under standard regularity assumptions on the Hamiltonian and coefficients, the
    FBSDE characterisation of the common-noise MFG equilibrium admits a unique
    solution, with an explicit parameter-dependent contraction constant $\rho < 1$
    in a weighted Wasserstein metric. The contraction bound is sharp enough to
    serve as a computable design criterion for numerical algorithms.

    \item \textbf{Numerical convergence} (Chapter 4, Theorem 4.1).
    A combined Euler--Maruyama/particle/SPDE scheme for the coupled FBSDE--Fokker--Planck
    system converges with weak Wasserstein error $O(\Delta t^{1/2} + \Delta x + M^{-1/2})$,
    where $\Delta t$ is the time step, $\Delta x$ the spatial mesh, and $M$ the
    particle count. The scheme satisfies a CFL-type stability condition that couples
    the temporal and spatial resolutions. Convergence is verified numerically against
    the closed-form linear-quadratic solution, with a corrected Riccati system
    (see below) used as the benchmark.

    \item \textbf{Particle methods and propagation of chaos} (Chapter 5).
    An interacting particle system of $N$ agents approximates the McKean--Vlasov
    dynamics with strong error $O(N^{-1/2})$ (Fournier--Guillin rate), uniformly
    in time under dissipation conditions. Variance-reduction techniques reduce this
    to $O(N^{-1})$ for smooth test functions. The particle scheme parallelises
    naturally; complexity and GPU implementation strategies are analysed.

    \item \textbf{Corrected linear-quadratic master equation} (Appendix A).
    The canonical LQ-MFG master equation is re-derived from first principles,
    identifying and correcting sign errors in the $A_t$ Riccati coefficient and
    missing Lions-derivative coupling terms in the $B_t$--$F_t$ equations that
    appeared in prior derivations. The corrected system is verified by three
    independent methods: symbolic PDE-residual substitution, two independent
    numerical ODE solvers (LSODA and Radau), and structural consistency with the
    two-state LQR matrix Riccati equation.

    \item \textbf{Lean 4 formal verification} (Chapter 6).
    Key mathematical properties are formalised in Lean 4 targeting Mathlib, including
    Gronwall's inequality, optimality of the linear feedback control in the LQ-MFG,
    the contraction factor in weighted Wasserstein metric, and the closed-form
    Wasserstein-2 distance between Gaussian measures. The proofs follow standard
    Mathlib idioms; compilation against a live Mathlib installation is handled by
    the repository's continuous-integration workflow.
\end{enumerate}

The honest limitations of the present work are documented throughout. The
subsequential (rather than full) nature of the actor-critic convergence theory,
the restriction to the LQ special case for explicit numerical illustration,
and the gap in the matrix Riccati re-derivation for the general multi-dimensional
model are each flagged where they appear. The Lean 4 proofs are reviewed for
mathematical correctness but should be rebuilt via \texttt{lake build} before
being relied on as machine-verified, as Mathlib's API evolves.

The work establishes a foundation for several natural extensions: jump-diffusion
common noise, mean-field control (social optimum), deep Galerkin and physics-informed
neural network solvers, and the Lata\l{}a--Oleszkiewicz inequality as a route from
subsequential to full convergence. These are discussed in Chapter 7.

\end{abstract}
